\documentclass[11pt]{article}
\usepackage[margin=1in]{geometry}
\usepackage{amsmath,amssymb,amsthm}
\usepackage{graphicx}
\usepackage[hidelinks]{hyperref}

\newtheorem*{theorem}{Theorem}
\newtheorem*{corollary}{Corollary}

\newcommand{\A}{\mathfrak A}
\newcommand{\C}{\mathcal C}
\newcommand{\ext}{\operatorname{ext}}

\title{A Literature-Implied Partial Answer for Complex Selfadjoint \(L_1\)-Preduals}
\author{Run \texttt{fa\_banach\_001}, lane 16}
\date{2026-06-19}

\begin{document}
\maketitle

\paragraph{Status.}
This is a \emph{literature-implied partial subcase}, not a claimed full solution
of the complex problem in arXiv:1105.3413.

\paragraph{Source question.}
Ludvik and Spurny prove their Theorem 1.4 for real \(L_1\)-preduals and write
on page 5 that the complex case remains open because it is not clear whether
Lusky's theorem that every real \(L_1\)-predual is complemented in a simplex
space remains true for complex \(L_1\)-preduals.

\begin{center}
\includegraphics[width=0.95\linewidth]{figures/open_problem_crop.png}
\end{center}

\paragraph{Implied subcase.}
The complex analogue of Theorem 1.4 holds for selfadjoint complex
\(L_1\)-preduals. In particular, by Hirsberg and Lazar's representation theorem,
it holds for every complex Lindenstrauss space whose unit ball has an extreme
point.

\begin{theorem}
Let \(X\) be a Choquet simplex and let
\[
E=A_{\mathbb C}(X):=A(X)+iA(X)
\]
be the complex Banach space of continuous complex-valued affine functions on
\(X\), with the supremum norm. Let \(f\in E^{**}\) be strongly affine on
\(B_{E^*}\) and suppose \(f\in \C_\alpha(B_{E^*})\) for some
\(\alpha\in[2,\omega_1)\). Then
\[
f\in
\begin{cases}
\A_{\alpha+1}(B_{E^*}), & 2\leq \alpha<\omega_0,\\
\A_{\alpha}(B_{E^*}), & \omega_0\leq \alpha<\omega_1.
\end{cases}
\]
If, in addition, the extreme boundary of the representing simplex \(X\) is
Lindelof and resolvable, then \(f\in\A_\alpha(B_{E^*})\).
\end{theorem}

\paragraph{Proof.}
Embed \(X\) into \(B_{E^*}\) by evaluations, \(\delta_x(a)=a(x)\). Put
\(F(x)=f(\delta_x)\). Since the evaluation map is continuous, \(F\) is of Baire
class \(\alpha\) on \(X\). Strong affinity of \(f\) on \(B_{E^*}\) implies
strong affinity of \(F\) on the simplex \(X\): if \(\mu\in M^1(X)\) has
barycenter \(x\), then the push-forward of \(\mu\) under \(y\mapsto\delta_y\)
has barycenter \(\delta_x\), hence
\[
F(x)=f(\delta_x)=\int_X f(\delta_y)\,d\mu(y)=\int_X F(y)\,d\mu(y).
\]

Apply Ludvik-Spurny's real simplex theorem, namely Proposition 7.1 in the proof
of their Theorem 1.4, to the real and imaginary parts of \(F\). This gives
\[
F\in
\begin{cases}
\A_{\alpha+1}(X;\mathbb C), & 2\leq \alpha<\omega_0,\\
\A_{\alpha}(X;\mathbb C), & \omega_0\leq \alpha<\omega_1,
\end{cases}
\]
and, under the Lindelof-resolvable boundary hypothesis, \(F\in\A_\alpha(X;\mathbb C)\).

It remains to lift the affine-class approximants from \(X\) back to \(B_{E^*}\).
Every \(\phi\in B_{E^*}\) is represented by a complex Radon measure \(\mu\) on
\(X\) with \(\|\mu\|\leq 1\), after extending \(\phi\) from \(E\subset C(X)\)
by Hahn-Banach and applying the Riesz representation theorem. If \(h\in
A_{\mathbb C}(X)\), define \(\widehat h(\phi)=\phi(h)=\int_X h\,d\mu\). If
\((h_n)\) is a bounded sequence of lower affine-class functions converging
pointwise to \(h\) on \(X\), dominated convergence gives
\[
\widehat h_n(\phi)=\int_X h_n\,d\mu \longrightarrow \int_X h\,d\mu
\]
for every \(\phi\in B_{E^*}\). Thus the lifting \(h\mapsto \widehat h\)
preserves all intrinsic affine classes generated by bounded pointwise sequential
limits.

Finally, the lifted function agrees with the original \(f\). For a representing
measure \(\mu\), decompose \(\mu\) into a linear combination of four positive
measures via the Jordan decompositions of its real and imaginary parts. After
normalizing the nonzero positive parts to probability measures, strong affinity
and complex linearity of \(f\in E^{**}\) give
\[
f(\phi)=\int_X F\,d\mu.
\]
Hence \(f\) is exactly the lift of \(F\), and it belongs to the asserted
intrinsic affine class on \(B_{E^*}\).

\begin{corollary}
If a complex Lindenstrauss space \(E\) has an extreme point in its unit ball,
then the complex analogue of Ludvik-Spurny's Theorem 1.4 holds for \(E\).
\end{corollary}

\paragraph{Justification.}
Hirsberg and Lazar proved that a complex Lindenstrauss space whose unit ball has
an extreme point is isometric to \(A_{\mathbb C}(X)\) for a Choquet simplex
\(X\). The theorem therefore applies.

\paragraph{Remaining open case.}
This does not settle the general complex \(L_1\)-predual case. Wulbert proved
that nonselfadjoint complex \(L_1\)-preduals exist, so the missing case is a
genuine residual part of the original question.

\paragraph{Search bounds.}
Cheap run indexes had no packet for arXiv:1105.3413 or this complex
\(L_1\)-predual question. Web searches for Lusky's theorem in complex
\(L_1\)-preduals and local corpus searches for ``complex Lindenstrauss'',
``selfadjoint'', and ``nonselfadjoint'' found the Hirsberg-Lazar and Wulbert
structure results but no full later settlement of the nonselfadjoint case.

\begin{thebibliography}{9}
\bibitem{LS}
P. Ludvik and J. Spurny,
\emph{Descriptive properties of elements of biduals of Banach spaces},
arXiv:1105.3413, 2011.

\bibitem{HL}
B. Hirsberg and A. J. Lazar,
\emph{Complex Lindenstrauss spaces with extreme points},
Transactions of the American Mathematical Society 186 (1973), 141--150.

\bibitem{W}
D. E. Wulbert,
\emph{Real structure in complex \(L_1\)-preduals},
Transactions of the American Mathematical Society 235 (1978), 165--181.
\end{thebibliography}

\end{document}
